\section{Resumen}

Esta tesis estudia el seguimiento de precios en un creador automático de mercado (AMM) sintético de producto constante y sin comisiones, frente a una referencia exógena de un mercado centralizado (CEX). Separa las consecuencias algebraicas de la regla de intercambio de los efectos que dependen del comportamiento de los agentes. En una piscina aislada, un arbitrajista limitado a una fracción c de la reserva de entrada puede reducir el error logarítmico fuera de la banda hasta en dos veces el logaritmo natural de uno más c por tick. Para un choque logarítmico de 0.20 y una banda de 0.003, el tiempo idealizado hasta la banda es de 20 ticks con c = 0.005, cuatro con c = 0.03 y dos con c = 0.10. Este cálculo explica la dirección del resultado de arbitraje configurado; la simulación mide desviaciones bajo ruido y otras reglas.

En 13,500 ejecuciones con 300 semillas Monte Carlo, el escenario de capacidad alta tiene un área media de exceso de error de 0.10, frente a 2.36 sin arbitraje. El resultado depende de la regla y la grilla de parámetros especificadas, y el signo del efecto está incorporado mecánicamente en esa regla. El disparador actual de retiro se activa en todas las corridas E3, normalmente en el primer tick. Una medida posterior y exploratoria del máximo error absoluto de seguimiento es mayor en E3 solo con factor de profundidad 0.5 y órdenes absolutas fijas (diferencias pareadas medias de 0.0104 a 0.0519 entre celdas de choque); el contraste es numéricamente cero con factores 1 y 2. El retiro también reduce el límite de arbitraje ligado a las reservas, así que este resultado secundario no aísla su efecto ni demuestra una amplificación general de liquidez.

El resultado de persistencia en E3 es específico del simulador: al ponderar por igual las nueve celdas del diseño, el 31.2 por ciento de las trayectorias recuperadas salió después de la banda de no arbitraje; con factor de profundidad 0.5, la proporción condicional fue 98.8 por ciento, mientras que con factores 1 y 2 no hubo salidas posteriores. Esto motiva registrar por separado el primer retorno y las salidas posteriores en futuros estudios de reproducción de eventos, pero no es una tasa de mercado real. El mapa reducido de dos estados mejora la persistencia en algunas condiciones, pero no supera a un AR(1) escalar. Un mapa sensible a la banda reduce cerca de 1.7 por ciento el error recursivo de reproducción, concentrado en una condición y usando la trayectoria futura realizada del CEX; no es un resultado de pronóstico en tiempo real. La contribución es un estudio de simulación reproducible y centrado en mecanismos que distingue resultados algebraicos de sensibilidades dependientes de reglas, y muestra dónde fallan sus modelos reducidos. La reproducción de eventos con datos de mercado sincronizados y costos realistas queda para trabajo futuro.

\textbf{Palabras clave:} creador automático de mercado; arbitraje; seguimiento de precios; retiro de liquidez; simulación basada en agentes; experimento Monte Carlo
